\documentclass[12pt,b5paper,twoside,openany]{extbook}

\usepackage{fancyhdr}
\usepackage{titlesec}		%for drop numbers & headers
\usepackage[Rejne]{fncychap}  % chapter style
\usepackage{geometry}	%for margins
\usepackage{multicol} %for two columns
\usepackage[toc]{multitoc} %two column table of contents
\usepackage{hyperref} %links

\title{Echoes of Eve}
\date{}
\author{}
\newgeometry{	inner=20mm,  % Margin closest to the binding
	outer=20mm,  % Margin on the outside edge (allows for journaling space)
	top=20mm,
	bottom=20mm} 
\begin{document}

\frontmatter
\maketitle
\tableofcontents

\chapter*{Preface}\addcontentsline{toc}{chapter}{Preface}
For almost a decade, I have been led to study what the Bible \textit{actually} says about women and to learn from the examples of the women in scripture. This booklet is provided as a resource for you, to make it easier for you to find people and verses in your own investigation. It contains a (nonexhuastive) list of women discussed in scripture and a few topics to study with Bible references  This list of women, topics, and corresponding verses is a product of my studies in the Word as well as the research of others, but it is not all encompassing. There are several challenges when approaching this area of study.

First, tradition informs our interpretation. Why does that make this topic difficult? Becuase the prevailing tradition in Christianity is patriarchal, which whether intended or unintended presents us with a view of women being secondary. Our first instinct when reading a name in the text is that it is male. Yet, there are many times when women are listed genealogies as mother or daughter (e.g., 1 Chronicles 4:3; 1 Chronicles 7:17-18). In addition, we have the challenges of language and culture--the names presented are both foreign in language and in time, thus names are not always obviously gendered the way modern names are. This instills in us a bias to assume all the names in genealogies belong to men unless explicitly stated otherwise. Yet, names such as Abihail appear equally masculine and feminine thorouhgout the text. So when no gender is explicitly assigned to the Abihail listed in Numbers 3:35, how do we know what gender to assume? Within this resource, I have listed people like Abihail with a footnote calling out the fact that there is doubt and confusion on the gender of the person, but it is also possible I have missed people entirely.

Another challenge lies in the number of unnamed women in scripture. These women are sometimes explicitly referenced without name, such as in Gensis 7:13 where we are told that the wives of Noah nad his 3 sons--4 unnamed women--were present on the ark. In this case, I have used the pattern "wife of <Husband's Name>" as the pattern and followed indexing format such that she would be listed alphabeticallys as "<Husband's Name>, wife of."  This comes with a secondary challenge for women who are discussed without a relative to reference, or multiple to reference. In those cases I tried to use the most common or logical descriptor. Thus, it you are searching for an unnamed woman, it may be a little more difficult to find her. Of note, there  are also implicitly references. For every genealogy that only lists one parent, there is a second parent that is implied even if they aren't mentioned. For example, Genesis 5:6-8 tells us that Seth had children without ever mentioning a woman. Yet, we know at least one woman existed for him to have children. I have not included these women in this compilation.

Finally, there is the question of what is relevant as a topic of study. There are many topics that pertain to both male and female, while others are more directly female (e.g., childbirth). The final pages of this book include topics generally directed toward women by either the Church or Society (often both) and various women and/or verses ofr deeper study on the topic. These are not the \textit{only} topics pertaining to women in the Bible nor are they topics that only pertain to women.  

I hope this resource helps you to strengthen your relationship with YHWH, encourages you to love and honor yourself and the women around you, and inspires you.


\mainmatter	
\chapter*{A}\addcontentsline{toc}{chapter}{A}

\begin{description}
	\item [Abel, the Wise Woman of] 2 Samuel 20
	\item [Abi] 2 Kings 18:2; 2 Chronicles 29:1
	\item [Abigail \#1]  \textit{wife of Nabal and later David}, 1 Samuel 25; 27:3; 30:1-5; 2 Samuel 3:3; 1 Chronicles 31:1
	\item [Abigail \#2] \textit{sister of David},  2 Smauel 17:25; 1 Chronicles 2:15-17
	\item [Abihail \#1] \textit{wife of Abishur}, 1 Chronicles 2:29
	\item [Abihail \#2] \textit{wife of Rehoboam}, 2 Chronicles 11:18
	\item [Abihail \#3\footnotemark{}] \textit{parent of Zuriel}, Numbers 3:35
	\footnotetext{the gender of this person is debated}
	\item [Abiram's Wife] Numbers 16:27
	\item [Abishag] 1 Kings 1-2
	\item [Abital] 2 Samuel 3:4; 1 Chronicles 3:3
	\item [Abraham's Daughter] Luke 13:10-17
	\item [Achsah] Joshua 15:13-19; Judges 1:12-15; 1 Chronicles 2:49
	\item [Adah \#1] \textit{wife of Lamech}, Genesis 4:19-24
	\item [Adah \#2] \textit{wife of Esau}, Genesis 28:8-9; 36:2
	\item [Adam \& Eve's Duaghters] Genesis 5:4
	\item [Adulterous Woman] John 8:1-11
	\item [Ahinoam \#1] \textit{wife of Saul}, 1 Samuel 14:50
	\item [Ahinoam \#2] \textit{wife of David}, 1 Samuel 25:43
	\item [Ahlai] 1 Chronicles 2:31
	\item [Aholibamah]Genesis 28:8-9; Genesis 36:2, 41; 1 Chronicles 1:52
	\item [Anah\footnotemark]  Genesis 36:2, 14, 18; 1 Chronicles 38, 40, 41
	\footnotetext{the gender of this person is debated}
	\item [Anna] Luke 2:36-38
	\item [Annointing Woman] Luke 7:36-39
	\item [Aseanth] Genesis 41:45, 50; 46:20
	\item [Ashtoreth] 1 Kings 11:5, 33; 23:13
	\item [Atarah] 1 Chronicles 2:26
	\item [Athaliah] 2 Kings 8:26; 11; 2 Chronicles 22; 23:12-15, 21; 24:7
	\item[Azubah \#1] \textit{mother of King Jehoshaphat}, 1 Kings 22:32
	\item[Azubah \#2] \textit{wife of Caleb},  1 Chronicles 2:18-19
\end{description}


\chapter*{B}\addcontentsline{toc}{chapter}{B}

\begin{description}
	\item [Baara] 1 Chronicles 8:8
	\item [Babylon, the Queen of] Daniel 5
	\item [Babylon, the Whore of] Revelation 17
	\item [BahurimBahurim] 2 Samuel 17:17-20
	\item [Barzillai, the Daughters of] Ezra 2:61; Nehemiah 7:63
	\item [Basemath \#1] \textit{daughter of Elon the Hittite}, Genesis 26:34; 36:2-4, 10
	\item [Basemath \#2] \textit{daughter of Ishmael}, Genesis 36:3-4, 10,13, 17
	\item [Basemath \#3] \textit{daughter of Solomon}, 1 Kings 4:15
	\item [Bathsheba] 2 Samuel 11; 12:1-25; 1 Kings 1:11-31; 2:13-22; Psalm 51
	\item [Berenice] Acts 25-26
	\item [Bilhah] Genesis 29:29; 30: 1-7; 35:22, 25; 37:2; 46:25; 1 Chronicles 7:13
	\item [Bithiah] \textit{or Bithia}, 1 Chronicles 4:18
	\item [Blind Man, mother of the] John 9:1-3
	\item [Bride of Christ] Ephesians 5:25-32; 2 Corinthians 11:2; Revelation 19:7-9\footnote{There are no verses in the Bible explicitly referencing a "Bride of Christ", verses listed here are how people put together the idea of a Bride of Christ}
\end{description}


\chapter*{C}\addcontentsline{toc}{chapter}{C}
\begin{description}
	\item [Cain, wife of] Genesis 4:17; 5:4
	\item [Canaanite Woman] Matthew 15:21-28; Mark 7:24-30
	\item [Candace] Acts 8:27
	\item [Careless Women] Isaiah 32:9-11
	\item [Chloe] 1 Corinthians 1:11
	\item[Claudia] 2 Timothy 4:21
	\item [Cozbi] Numbers 25
\end{description}


\chapter*{D}\addcontentsline{toc}{chapter}{D}
\begin{description}
	\item [Damaris] Acts 17:32-34
	\item [Dathan, wife of] Numbers 16:27-33
	\item [Deborah \#1] \textit{nurse of Rebekah}, Genesis 24:59; 35:8
	\item [Deboarh \#2] \textit{judge of Israel}, Judges 4-5
	\item [Delilah] Judges 16
	\item [Dinah] Genesis 30:21; 34; 46:15
	\item [Dorcas] \textit{also known as Tabitha}, Acts 9:36-42
	\item [Drusilla] Acts 24:24
\end{description}


\chapter*{E}\addcontentsline{toc}{chapter}{E}
\begin{description}
	\item [Eglah] 2 Samuel 3:5; 1 Chronicles 3:3
	\item [Elisheba] Exodus 6:23
	\item [Elizabeth] Luke 1
	\item [Ensnaring Women]	Ecclesiastes 7:26-29
	\item [Ephah] 1 Chronicles 2:46
	\item [Ephrath] \textit{or Ephrathah}; 1 Chronicles 2:19, 24, 50; 4:4
	\item [Esther] \textit{Hebrew name Hassadah}; The book of Esther
	\item [Eunice] 2 Timothy 1:5
	\item [Euodia] Philippians 4:2
	\item [Eve] Genesis 1:26-31; 2-5
	\item [Ezekiel, wife of] Ezekiel 24:15-24
\end{description}


\chapter*{F}\addcontentsline{toc}{chapter}{F}
\begin{description}
	\item [False Prophets] Ezekiel 13:17-23
	\item [Foreign Wives\footnotemark{}] Nehemiah 13:23-28
	\footnotetext{description taken from the NASB, the KJV uses "strange wives"}
	\item [Fortune Teller] Acts 16:16-21
\end{description}


\chapter*{G}\addcontentsline{toc}{chapter}{G}
\begin{description}
	\item [Gideon, wives \& concubine of] Judges 8:30-31
	\item [Gomer] Hosea 1; 3
\end{description}


\chapter*{H}\addcontentsline{toc}{chapter}{H}
\begin{description}
	\item [Hagar] Genesis 16; Genesis 21:9-21; 25:12; Galatians 4:24-25
	\item [Haggith] 2 Samuel 3:4; 1 Kings 1:5, 11; 2:13; 1 Chronicles 3:2
	\item [Ham, wife of] Genesis 7:13
	\item [Hammolekheth] \textit{or Hammolecheth}, 1 Chronicles 7:17-18
	\item [Hamutal] 2 Kigns 23:31; 24:18; Jeremiah 52:1
	\item [Hannah] 1 Samuel 1; 2:1-10
	\item [Harlots, the two] 1 Kings 3:16-28
	\item [Hazzelelponi] 1 Chronicles 4:3
	\item [Helah] 1 Chronicles 4:5, 7
	\item [Hephziba] 2 Kings 21:1
	\item [Herodias] Matthew 14:3-8; Mark 6:17-28; Luke 3:19-20
	\item [Herodias, daughter of\footnotemark{}]
	\footnotetext{Sources suggest her name might have been Salome; since there is no Biblical evidence for this, she is NOT listed under Salome}  Matthew 14:6-8; Mark 6:17-28
	\item [Hodesh] 1 Chronicles 8:8-9
	\item [Hodiah, wife of] 1 Chronicles 4:19
	\item [Hoglah] Numbers 26:33; 27:1-8; 36:11; Joshua 17:3-4
	\item [Huldah] 2 Kings 22:14-20; 2 Chronicles 34:22-28
	\item [Hushim] 1 Chronicles 8:8,11
\end{description}


\chapter*{I}\addcontentsline{toc}{chapter}{I}
\begin{description}
	\item [Isaiah, wife of\footnotemark{}] 	\footnotetext{Note, the text does not confirm Isaiah was married to tihs woman; this interpretation is likely because as a prophet we expect him to only have sex within marriage.}Isaiah 7:14-16\footnote{this is generally considered a dual prophecy, with the first fulfillment in Isaiah's son, thus "the virgin" could also refer to this mother of Isaiah's child}; 8:3

	
	\item [Iscah\protect\footnotemark{}] Genesis 11:29
	\footnotetext{The referenced text implies Iscah is a woman but does not explicitly tell us this; upon researching I found lexicons show this Hebrew name (Strong's H3252) to be defined as the "daughter of Haran" or the sister of Lot and Milcah.}
\end{description}


\chapter*{J}\addcontentsline{toc}{chapter}{J}
\begin{description}
	\item [Jabez, mother of] 1 Chronicles 4:9
	\item [Jael] Judges 4-5
	\item [Japheth, wife of] Genesis 7:13
	\item [Jarius, duaghter of] Matthew 9:18-19, 23-26; Mark 5:21-24, 35-43; Luke 8:40-42, 49-56
	\item [Jecoliah] \textit{or Jecholiah}, 2 Kings 15:1-2
	\item [Jedidah] 2 Kings 22:1
	\item [Jehoaddan] \textit{or Jehoaddin}, 2 Kings 14:1-2; 2 Chronicles 25:1
	\item [Jehosheba] 2 Kings 11:1-2
	\item [Jehudijah\protect\footnotemark{}] 1 Chronicles 4:17-18
	\footnotetext{this "name" shows up in the KJB but is more than likely a descriptor; as such other translations refer to the "Jewish Wife"}
	\item [Jemima] Job 42:13-15
	\item [Jephthah, daughter of] Juges 11
	\item [Jephthah, mother of] Judges 11:1-2
	\item [Jerioth] 1 Chronicles 2:18
	\item [Jeroboam, wife of] 1 Kings 14:1-18
	\item [Jerusha] 2 Kings 15:33; 2 Chronicles 27:1
	\item [Jesus, sisters of\protect\footnotemark] Matthew 13:55-56
	 \footnotetext{Jesus is an English name that stands in for the Hebrew name Yeshua; this entry is duplicated under Y under the name Yeshua}
	\item [Jezebel \#1] 1 Kings 16:31; 1 Kings 18:2-3, 13, 17-19; 19:1-3; 21; 2 Kings 9
	\item [Jezebel \#2] Revelation 2:20
	\item [Joanna] Luke 8:1-3; 24:1-10
	\item [Job, sisters of] Job 42:11
	\item [Job, wife of] Job 2:9; 19:17; 31:10
	\item [Job, 3 daughters of] Job 1:2, 13, 18
	\item [Jochebed] Exodus 6:20; Numbers 26:59
	\item [Judah, wife of] Genesis 38:1-5
	\item [Judith] Genesis 26:34
	\item [Julia] Romans 16:15
	\item [Junia] Romans 16:7
\end{description}


\chapter*{K}\addcontentsline{toc}{chapter}{K}
\begin{description}
	\item [Keren-Happuch] Job 42:13-15
	\item [Keturah] Genesis 25:1-4; 1 Chronicles 1:32-33
	\item [Keziah] Job 42:13-15
\end{description}


\chapter*{L}\addcontentsline{toc}{chapter}{L}
\begin{description}
	\item [Leah] Genesis 29:16-17, 21-35; 30:9-21; 31; 33:1-7; 34:1; 35:23-26; 46:15-18; 49:21; Ruth 4:11
	\item [Lemuel, mother of] Proverbs 31
	\item [Lo-Ruhamah] Hosea 1:6-8
	\item [Lois] 2 Tmothy 1:3
	\item [Lost Coin Woman] Luke 15:8-10
	\item [Lot, daughters of] Genesis 19
	\item [Lot, wife of] Genesis 19:1-26
	\item [Lydia] Acts 16:14-15, 40
\end{description}


\chapter*{M}\addcontentsline{toc}{chapter}{M}
\begin{description}
	\item [Maacah \#1]  \textit{wife of David}, 2 Samuel 3:3; 1 Chronicles 3:2
	\item [Maacah \#2] \textit{wife of Rehoboam} 1 Kings 15:2,10,13; 2 Chronicles 11:18, 21-22, 43; 2 Chronicles 15:16
	\item [Maacah \#3]  \textit{concubine of Caleb} 1 Chronicles 2:48
	\item [Maacah \#4]  \textit{wife of Machir} 1 Chronicles 7:15-16
	\item [Maacah \#5]  \textit{mothter of Gibeon} 1 Chronicles 8:29; 9:35
	\item [Maacah \#6]  \textit{parent of Hanan},\footnote{the gender of this person is debated} 1 Chronicles 11:43
	\item [Mahalath \#1] \textit{wife of Esau}, Genesis 28:9	
	\item [Mahalath \#2] \textit{wife of Rehoboam}, 2 Chronicles 11:18
	\item [Mahlah \#1] \textit{daughter of Zolephehad}, Numbers 26:33; 27:1-8; 36:11; Joshua 17:3-4
	\item [Mahlah \#2] \textit{child of Hammoleketh}\footnote{the gender of this person is debated}, 1 Chronicles 7:18
	\item [Martha] Luke 10:38-42; John 11:1-46; 12:1-2
	\item [Mary Magalene] Matthew 27:55-61; 28:1-10; Mark 15:40-47; Mark 16:1-11; Luke 8:1-3; 24:1-11; John 19:25; John 20:1-2, 11-18
	\item [Mary, mother of James \& Joses\footnotemark{}] Matthew 27:55-61; 28:1-10; Mark 15:40-47; 16:1-8; Luke 24: 1-11
	\footnotetext{debated whether this Mary is also Mary mother of Jesus/Yeshua}
	\item [Mary, mother of Jesus/Yeshua] Matthew 1:16-25; 2:11, 13-14, 19-23; 13:55; Mark 6:3; Luke 1:26-56; 2; John 19:25-27; Acts 1:14
	\item [Mary, mother of John Mark] Acts 12:12
	\item [Mary, of Rome] Romans 16:6
	\item [Mary, sister of Martha] Luke 10:38-42; John 11:1-46; 12:1-8
	\item [Mary, wife of Clopas] John 19:25
	\item [Matred] Genesis 36:39; 1 Chronicles 1:50
	\item [Mehetabel] Genesis 36:39; 1 Chronicles 1:50
	\item [Mephibosheth, nurse of] 2 Samuel 4:4
	\item [Merab] 1 Samuel 14:49; 18:17-19; 2 Samuel 21:8
	\item [Meshullemeth] 2 Kings 21:19
	\item [Micah, mother of] Judges 17:1-4
	\item [Michal] 1 Samuel 14:49; 18:20-28; 19:11-17; 25:44; 2 Samuel 3:13-16; 6:16, 20-23; 1 Chronicles 15:29
	\item [Milcah \#1] \textit{wife of Nahor}, Genesis 11:29; 22:20-23; 24:15, 24, 47
	\item [Milcah \#2] \textit{daughter of Zolephehad}, Numbers 26:33; 27; 36:11; Johsua 17:3
	\item [Miriam \#1] Exodus 2:1-10; 15:20-21; Numbers 12; 20:1; 26:59; Deuteronomy 24:9; 1 Chronicles 6:3; Micah 6:4
	\item [Miriam \#2] 1 Chronicles 4:17
	\item [Mothers who ate sons] 2 Kings 6:24-33
\end{description}


\chapter*{N}\addcontentsline{toc}{chapter}{N}
\begin{description}
	\item [Naamah \#1] \textit{daughter of Lamech and Zillah}, Genesis 4:22
	\item [Naaham \#2]\textit{mother of Rhoboam}, 1 Kings 14:21, 31; 2 Chronicles 12:13
	\item [Naaman, wife of] 2 Kings 5:2
	\item [Naaman, servant of wife of] 2 Kings 5:2
	\item [Naarah] 1 Chronicles 4:5-6
	\item [Nahash\footnotemark{}] 2 Samuel 17:25
	\footnotetext{the gender of this person is debated}
	\item [Namoi] Book of Ruth
	\item [Nehushta] 2 Kings 24:8
	\item [Nereus, sister of] Romans 16:15
	\item [Noadiah] Nehemiah 6:14
	\item [Noah] Numbers 26:33; 27:1-11; 36:11
	\item [Noah, wife of] Genesis 6-9
\end{description}

\chapter*{O}\addcontentsline{toc}{chapter}{O}
\begin{description}
	\item [Orpah] Ruth 1
\end{description}

\chapter*{P}\addcontentsline{toc}{chapter}{P}
\begin{description}
	\item [Paul, sister of] Acts 23:16
	\item [Peninnah] 1 Samuel 1:1-7
	\item [Persis] Romams 16:12
	\item [Peter, mother-in-law of] Matthew 8:14-17
	\item [Peter, wife of]  Matthew 8:14-17
	\item [Pharaoh, daughter of] Exodus 1-3
	\item [Phillip, daughters of] Acts 21:8-9
	\item [Phinehas, wife of] 1 Samuel 4:19-22
	\item [Phoebe] Romans 16:1-2
	\item [Pilate, wife of] Matthew 27:	19
	\item [Potiphar, wife of] Genesis 39:7-19
	\item [Preflood daughters\footnotemark{}] Genesis 5
	\footnotetext{All the sons of Adam are said to have daughters with no other defining information such as name or number; instead of creating an entry for each son, where only the verse references change, I am grouping them together in this entry.}
	\item [Priscilla] Acts 18:2, 18, 26; Romans 16:3-5
	\item [Puah] Exodus 1:15-22
\end{description}

%\chapter*{Q}
%\begin{description}
%\end{description}

\chapter*{R}
\begin{description}\addcontentsline{toc}{chapter}{R}
	\item [Rachel] Genesis 29-31; 33:1-7; 35:16-20, 24, 25; 46:19, 22, 25; 48:7; Ruth 4:11; 1 Samuel 10:2; Jeremiah 31:15; Matthew 2:18
	\item [Rahab] Joshua 2:1-22; 6:17, 23, 25; Matthew 1:5; Hebrews 11:31; James 2:25
	\item [Rahab, sisters of] Joshua 2:13; 6:23
	\item [Rebekah] Genessis 22:20-23; 24; 25:20-24; 26:6-11, 34-35; 27:5-17, 42-46; 28:5; 29:12; 35:8; 49:31; Romans 9:10
	\item [Rebekah, mother of] Genesis 22:28, 55
	\item [Reuel, daughters of] Exodus 2:16-22
	\item [Reumah] Genesis 22:23-24
	\item [Rhoda] Acts 12:13
	\item [Rizpah] 2 Samuel 3:7; 21:10-11
	\item [Rufus, mother of] Romans 16:13\footnote{Mark 15:21 is of interest as well, presumably this is the Rufus Paul is greeting thus making her the wife of Simon of Cyrene}
	\item [Ruth] the Book of Ruth; Matthew 1:5
\end{description}

\chapter*{S}\addcontentsline{toc}{chapter}{S}
\begin{description}
	\item [Salome\protect\footnotemark] Mark 15:33-41; 16:1-8
	\footnotetext{Some people believe Salome is the name of Zebedee's wife, and thus the mother of James and John. This would make the passage in Matthew 20:20-24 about her as well}
	\item [Samaritan Woman]  John 4:1-45
	\item [Samson, mother of] Judges 13-14
	\item [Samson,wife of] Juges 14-15
	\item [Sapphira] Acts 5:1-11
	\item [Sarah] \textit{or Sarai}, Genesis 11:29-31; 12:5-20; 16; 17:15-22; 18:1-15; 20; 21:1-14; 23; 24: 36, 67; 25:10, 12; 49:31; Isaiah 51:2; Romans 4:19; 9:9; Hebrews 11:11; 1 Peter 3:6
	\item [Serah] Genesis 46:17; Numbers 26:4; 1 Chronicles 7:30
	\item [Servants of the High Preist] Matthew 26:69-72; Mark 14:66-69: Luke 22:56-58
	\item [Servants of the Meeting] Exodus 38:8
	\item [Shallum, daughters of] Nehemiah 3:12
	\item [Sheba, Queen of] 1 Kings 10:1-13; 2 Chronices 9:1-12
	\item [Sheerah] 1 Chronicles 7:24
	\item [Shelomit \#1] \textit{or Shelomith; mother of the blasphemer}, Leviticus 24:11
	\item [Shelomit \#2] \textit{or Shelomith;daughter of Zerubbabel}, 1 Chronicles 3:19
	\item [Shem, wife of] Genesis 7:13
	\item [Shimeath] 2 Kings 12:21; 2 Chronicles 24:26
	\item [Shimrith] 2 Kings 12:21; 2 Chronicles 24:26
	\item [Shiphrah] Exodus 1:15-22
	\item [Shua, daughter of] Genesis 38:2-5, 12
	\item [Shunamite Woman] 2 Kings 4:8-37; 8:1-6
	\item [Solomon, wives and concubines of] 1 Kings 11:1-8
	\item [Strange Wives\footnotemark{}] Nehemiah 13:23-28
	\footnotetext{description taken from the KJV, other translations may use "foreign wives"}
	\item [Susanna] Luke 8:1-3
	\item [Syntyche] Philippians 4:2
	\item [Syrophoenician Woman] Mark 7:24-30
\end{description}

\chapter*{T}\addcontentsline{toc}{chapter}{T}
\begin{description}
	\item [Tahpenes] 1 Kings 11:19
	\item [Tahpenes, sister of] 1 Kings 11:20
	\item [Tamar \#1] \textit{matriarch of the Tribe of Judah}, Genesis 38; Ruth 4:12; 1 Chronicles 2:4; Matthew 1:3
	\item [Tamar \#2] \textit{daughter of David}, 2 Samuel 13; 1 Chronicles 3:9
	\item [Tamer \#3] \textit{daughter of Absalom}, 2 Samuel 14:27
	\item [Taphath] 1 Kings 4:11
	\item [Tekoa, Woman of] 2 Samuel 14:1-20
	\item [Thebez, Woman of] Judges 9:50-56
	\item [Timna] Genesis 36:12, 22; 1 Chronicles 1:39
	\item [Tirzah] Numbers 26:33; 27:1-11; 36:11; Judges 17:3
	\item [Titus 2 Woman] Titus 2:3-5
	\item [Tryphena] Romans 16:12
	\item [Tryphos] Romans 16:12
	
\end{description}

\chapter*{U}\addcontentsline{toc}{chapter}{U}
\begin{description}
	\item [Unnamed Concubine, the] Judges 19-21
\end{description}

\chapter*{V}\addcontentsline{toc}{chapter}{V}
\begin{description}
	\item [Vashti] Esther 1
	\item [Virgins, the Ten] Matthew 25:1-13
\end{description}

\chapter*{W}\addcontentsline{toc}{chapter}{W}
\begin{description}
	\item [Wicked Women] Jeremiah 44:15-30; Amos 4
	\item [Widow \#1] \textit{widow with a debt}, 2 Kings 4:1-7
	\item [Widow \#2] \textit{widow seeking justice}, Luke 18:1-8
	\item [Widow \#3] \textit{widow giving offering}, Mark 12:41-44; Luke 21:1-4
	\item [Widow of Nain] Luke 7:11-17
	\item [Widow Zarephath] 1 Kings 17:8-24; Luke 4:25-26
	\item [Widows] 1 Timothy 5
	\item [Wise Women] Exodus 35:25-26
	\item [Witch of En Dor] 1 Samuel 28:7-25
	\item [Woman from Song of Solomon] The book of Song of Solomon
\end{description}

%\chapter*{X}
%\begin{description}
%
%\end{description}

\chapter*{Y}\addcontentsline{toc}{chapter}{Y}
\begin{description}
	\item [Yeshua, sisters of\protect\footnotemark] Matthew 13:55-56
	\footnotetext{Yeshua is the Hebrew name for Jesus; this entry is duplicated under J for those more familiar with the name Jesus}
\end{description}

\chapter*{Z}\addcontentsline{toc}{chapter}{Z}
\begin{description}
	\item [Zebedee, wife of] Matthew 20:20-23; 27:45-56
	\item [Zebidah] 2 Kings 23:36
	\item [Zeresh] Esther 5:9-14; 6:13-14
	\item [Zeruiah] 1 Samuel 26:6; 2 Samuel 2:13, 18; 3:39; 8:16; 14:1; 16:9-10; 17:25; 18:2; 19:21-22; 21:17; 23:18,37; 1 Kings 1:7; 2:5, 22; 1 Chronicles 2:16; 11:6, 39; 18:12, 15; 26:28; 27:24 
	\item [Zibiah] 2 Kings 12:1; 2 Chronicles 24:1
	\item [Zillah] Genesis 4:19-24
	\item [Zilpah] Genesis 29:24; 30:9-12; 35:26; 37:2; 46:18
	\item [Zion, daughters of] Isaiah 3:16-26
	\item [Zipporah] Exodus 2:16-22; 4:24-26; 18:2-6; Numbers 12\footnote{it is debated whether this refers to Zipporah or a second wife}
\end{description}

\chapter *{Topics}\addcontentsline{toc}{chapter}{Topics}
\vspace{-.75in}
\begin{description}
	\item [The Creation and Purpose of Woman] Study \textit{Eve}, do a word study on "helpmate," and the following verses: Genesis 1-3 (See \href{https://www.psalmstogod.com/2023/01/the-creation-and-purpose-of-woman.html}{Creation and Purpose of Woman} for my study on the topic.\footnote{https://www.psalmstogod.com/2023/01/the-creation-and-purpose-of-woman.html})
	\item [Barrenness/Fertility] Study \textit{Sarah}, \textit{Rebekah}, \textit{Rachel}, \textit{Hannah}, \textit{Samson's Mother}, \textit{Elizabeth} (See \href{https://www.psalmstogod.com/2023/12/women-barrenness-in-bible.html}{Women \& Barrenness in the Bible} for my study on this.\footnote{https://www.psalmstogod.com/2023/12/women-barrenness-in-bible.html})
	\item [Rape/Sexual Assault] Study \textit{Dinah}, \textit{Tamar \#2}, \textit{the Unnamed Concubine}, \textit{the Wife of Potiphar} and the following verses: Deuteronomy 22:25-29; Galatians 5:19-21; Ephesians 5:28-29; Isaiah 61:7–8 (I have done a few \href{https://www.psalmstogod.com/search/label/Sexual%20Assault}{studies} that touch on this topic but not a single all encompassing study.\footnote{https://www.psalmstogod.com/search/label/Sexual\%20Assault})
	\item [Submission/Leadership/Teaching] Study \textit{Deborah \#2}, \textit{Huldah}, \textit{Abigail \#1}, \textit{Priscilla}, and the following verses: Ephesians 5:21-33; 1 Timothy 2:11-15; Titus 2:3-4
	\item [Singleness] Study \textit{Philip's daughters} and the following verses: 1 Corinthians 7; Ephesisans 5:22-33; 1 John 4:19; Romans 5:8; Matthew 7:11; Daniel 3 (See \href{https://www.psalmstogod.com/2025/10/3-passages-to-help-you-in-singleness.html}{3 Passages to Help You in Singleness} for my study on this.\footnote{https://www.psalmstogod.com/2025/10/3-passages-to-help-you-in-singleness.html})
	\item [Modesty] Study the following verses: 1 Timothy 2:9-10; Matthew 5:20; 18:9; 1 Corinthians 10:13; Genesis 2:25; 3:6-11, 21 (See \href{https://www.youtube.com/watch?v=nGw_gcEAYJc}{The Question of Modesty} for my study on this.\footnote{https://www.youtube.com/watch?v=nGw\_gcEAYJc})
\end{description}

\end{document}
